\chapter{Wyrms, Drakes, and Dragon-Kin}
\chaptermark{Dragons and Drakes}

Dragons are not beasts. Drakes are not dragons. The first are epochs with teeth---thinking sovereignties whose slumbers shape kingdoms. The second are animals: clever, dangerous, and sometimes supernatural, but animals nonetheless. Between them stand the dragon-kin, the mortal creatures who have traded breath, blood, or oath for a share of draconic fire.

This chapter is not a complete bestiary. It is a field guide to the three kinds of scaled things that matter most in a Fate's Edge campaign: the wise, the wild, and the sworn.

%====================================================
\section{Dragons --- The Thinking Calamity}
%====================================================

A dragon is a force of narrative gravity. When one wakes, borders redraw themselves. When one dies, the weather remembers. They do not hunt for food alone; they hunt for tribute, for stories, for debts that nations have forgotten.

Use the following template when an elder dragon must be statted. Treat the dragon as a setting choice first and a combatant second.

\subsection*{Elder Dragon Template}

\noindent
\textbf{Tier:} IV--V \hfill
\textbf{Harm:} 6--8 \hfill
\textbf{Armor:} 3--4

\medskip
\noindent
\textbf{Tags:} \textsc{Ward}, \textsc{Flame}, \textsc{Flier}, \textsc{Tyrant}, \textsc{Ancient}

\medskip
\noindent
\textbf{Dragon Moves} (choose two or three)
\begin{itemize}
  \item \textbf{Reposition the World.} Stairs invert, Near becomes Far, archways swallow pursuit.
  \item \textbf{Demand the Tithe.} The lair asks a price---blood, vow, name, or memory.
  \item \textbf{Answer with Weather.} The dragon emotes through storm, flood, quake, or frost for one beat.
  \item \textbf{Arouse the Hoard.} Treasures animate to defend the dragon's history.
  \item \textbf{Close the Eye.} Vision narrows; ranged actions become \emph{Desperate} until the characters change position.
\end{itemize}

\medskip
\noindent
\textbf{Waking Violation:} Name the one act that wakes the dragon immediately. Common violations include broken pacts, stolen heart-pieces, fire-rites in frozen lands, or a king who forgets fear.

\subsection*{Named Elders at a Glance}

\begin{tabular}{p{1.6in}p{2.2in}p{2.4in}}
  \textbf{Name} & \textbf{Domain} & \textbf{Core Want} \\
  Vyrmja the Winter Coil & Glacier cathedral of singing ice & The truth of every lineage \\
  Azghal of the Red Vault & Ruined city crown of broken oaths & Pride taken from another, as tribute \\
  Tyrgoth the Thunder-Eater & Sky-needles of broken storm-altars & One honest blow struck against his scales \\
  Sir Cadmorrant the Gilded & Valley of ruined heralds and gilded bones & A lineage he has never heard spoken
\end{tabular}

%====================================================
\section{Drakes --- The Animal Fire}
%====================================================

Drakes are the living barometer of frontier fate. They are not puzzles; they are omens. Clever play, bribes, and geasa matter more than raw harm. A drake does not scheme like a dragon; it reacts like a beast---territorial, hungry, and superstitious.

The table below gives six common drakes. Use them as wandering threats, lair guardians, or bad omens.

\begin{tabular}{p{1.3in}p{0.5in}p{0.6in}p{1.5in}p{1.8in}p{1.2in}}
  \textbf{Drake} & \textbf{Tier} & \textbf{Harm} & \textbf{Tags} & \textbf{Signature Move} & \textbf{Weakness or Trophy} \\
  Emberdrake & I & 1 & Ignite, Fast & Scorch the ground; flare to blind & Milk, ash, or sweet-smoke pacifies \\
  Mirewing Basilisk-Drake & II & 2 & Bind, Loping & Fix a victim with a muddy gaze & Polished mirrors; sudden bright clatter \\
  Barrow Serpent & II & 2 & Fear, Burrow & Suck warmth from the air & Funeral bells or true names of the buried \\
  Sky-Needle Wyvern & II--III & 2--3 & Wind, Fly & Snatch and drop from above & Nets, whistling cords, anchored lines \\
  Lampwyrm Archivist & I--II & 1 & Truth, Glide & Illuminate lies in its light & Candle snuffed respectfully; gifted annotation \\
  Frostling Wyrm & III & 3 & Ward, Cold, Slow & Breathe hoarfrost across an area & Resonant bells; shared fire and stories
\end{tabular}

\subsection*{Running a Drake}

When the party meets a drake, ask three questions:

\begin{enumerate}
  \item What territory does it claim?
  \item What has disturbed it?
  \item What does it want that is not meat?
\end{enumerate}

A drake can often be bought off with food, a shiny object, a draconic geis, or a show of submission. Killing one may draw the attention of a true dragon.

%====================================================
\section{Dragon-Kin --- The Mortal Flame}
%====================================================

Dragon-kin are mortals who carry the wyrm's shadow. Some gain power through pacts; others are shaped by worship, service, or bloodline curse.

\subsection*{Runekeepers}

A Runekeeper bargains with an ancient wyrm for magical power. The wyrm becomes a mystic patron. The price is always more than gold.

\noindent
\textbf{Pact Price (Obligation):}
\begin{itemize}
  \item Protect a treasure, name, land, or secret.
  \item Enforce a vow sworn by others.
  \item Extend the dragon's influence through fear, tribute, or renown.
  \item Once per season, deliver a worthy offering.
\end{itemize}

\noindent
\textbf{Refusing a demand} marks $+2$ Obligation. \textbf{Betraying a law} awakens the wyrm.

Sample Rites:
\begin{itemize}
  \item \textbf{Spark of the First Flame.} Ignite a melee weapon for one scene.
  \item \textbf{Scaled Skin.} Gain Armor 1 against mundane attacks.
  \item \textbf{Dragon's Gaze.} Force a creature to obey a short, simple command.
  \item \textbf{Crown of Scales.} Gain Armor 2 and immunity to fire for one scene.
  \item \textbf{Molten Breath.} Exhale a cone of flame for area harm.
\end{itemize}

\subsection*{Dragon Cults}

A dragon cult is built from three pillars:
\begin{enumerate}
  \item \textbf{Myth Engine.} What story powers the cult?
  \item \textbf{Scarce Rite.} What can only initiates do?
  \item \textbf{Territorial Claim.} Where does its influence hold?
\end{enumerate}

Cults move slowly. Track three timers:
\begin{itemize}
  \item \textbf{Infiltration [6].} Converts militia, merchants, clergy, children.
  \item \textbf{Manifestation [8].} Tries to physically call, awaken, or anchor a dragon.
  \item \textbf{Rupture [10].} Success changes local law, season, or geography.
\end{itemize}

\noindent
\textbf{Rites (sample):}
\begin{itemize}
  \item \textbf{Ember-Send.} Burn a sealed scrap; smoke carries a whispered message.
  \item \textbf{Scaled Veil.} Initiates appear as favored kin, inspiring dread.
  \item \textbf{Hoard-Calling.} Metal leaps to form a nest-mound.
  \item \textbf{Scale-Dawn.} A spectral draconic shape blots the sun; panic spreads.
\end{itemize}

\subsection*{Terrestrial Allegiances}

Mortal patrons use dragons the way dragons use kingdoms. Build a terrestrial allegiance quickly:

\begin{enumerate}
  \item \textbf{Who?} A guildmaster, commodore, magistrate, archivist, or duchess.
  \item \textbf{What do they want?} Bind the dragon, monopolize a tithe, recover a stolen oath, or weaponize dragon-weather.
  \item \textbf{Boon.} Choose one: Seal of Passage, Leveraged Favor, Licensed Rite, Logistics Surge, Heraldry of Safe-Conduct.
  \item \textbf{Obligation.} Pay a tithe, carry their mark, share first pick, or do not slay the dragon without their assent.
  \item \textbf{Escalation Track.} Notice $\rightarrow$ Audit $\rightarrow$ Sanction $\rightarrow$ Seizure.
\end{enumerate}

\subsection*{Knightly Orders and Geasa}

Knightly orders are political weathervanes as much as warriors. A geis is the sharp edge of honor.

\noindent
\textbf{Swearing a Geis:}
\begin{enumerate}
  \item Propose concrete terms with a trigger and a limit.
  \item Choose a seal: blood, name, relic, or witness.
  \item Mark Obligation segments per magnitude.
  \item Gain a situational tag or bonus die in the vowed context.
  \item Breaking the geis triggers Backlash and ticks all segments at once.
\end{enumerate}

\noindent
\textbf{Magnitude Table:}
\begin{tabular}{lll}
  \textbf{Scope} & \textbf{Boon} & \textbf{Obligation} \\
  Duel/Scene & $+1$ die or one tag & 2 segments \\
  Quest/Arc & $+1$ die and one tag & 4 segments \\
  Season/War & Two tags and a special clause & 6 segments
\end{tabular}

%====================================================
\section{Lairs and Hoards at a Glance}
%====================================================

A dragon's domain is alive. A dragon's hoard is a biography. Use this quick build procedure when the party stumbles into scaled territory.

\subsection*{Quick Lair Build}

\begin{enumerate}
  \item Choose a domain: volcanic caldera, storm fortress, glacier cathedral, sunken metropolis, ruined city crown, etc.
  \item Pick one lair heart: Hoard Chamber, Throne of Element, Court of Oaths, Egg Crypt, Observatory of Scars, Ossuary Forge.
  \item Choose two approaches/hazards: knife-ledge path, ashfall, mirror-pool maze, siren vents, glass rain, avalanche memory.
  \item Pick one or two denizens: Drake-Knights, Hoard-Wyrds, Sky Serpents, Herald Gargoyles, Ember Imps.
  \item State the tithe and the violation.
  \item Choose two lair moves.
\end{enumerate}

\subsection*{Quick Hoard Build}

\begin{enumerate}
  \item Choose one item from each stratum: Coin, Art, Name/Oath, Cursed Complication.
  \item Assign one or two tags as properties.
  \item State a Hoard Law: who may touch what, and what happens when it is broken.
\end{enumerate}

\noindent
\textbf{Spending hoard treasure} does not buy grain. It buys exceptions: a favor, a seal of transit, a year of silence, a cleared name, or a writ of passage. Each time treasure is spent this way, tick \textbf{Dragon's Attention [6]}. At six ticks, the wyrm knows where its wealth went.

%====================================================
\section{Adventure Seeds}
%====================================================

\begin{itemize}
  \item \textbf{The Sleeping Crown.} A dragon older than the capital lies beneath the city's foundation. The wards fail one by one. The dragon does not want vengeance; it wants a throne rebuilt.
  \item \textbf{Ash Above the Orchard.} A village prospers because a young scarlet drake nests in the orchard caves. Then an elder dragon comes to reclaim its ``child.''
  \item \textbf{A Tax of Wings.} Every spring, black-scaled tithe-collectors darken the sky. The tyrant-dragon is gone, but its brood still comes. The nobility offered the pact willingly, long ago.
  \item \textbf{The Knight With No Shadow.} A famous commander never casts a shadow because his shadow is a dragon bound in human form. The dragon does not want freedom; it wants the knight's body.
\end{itemize}

%====================================================
\section*{GM Reminder}
%====================================================

\noindent
\textit{Dragons are smart. Drakes are animals. Dragon-kin are mortals who have chosen to stand in the fire. Keep the distinction sharp, and every scaled encounter will feel like a different kind of story.}
